\documentclass[11pt]{article}

% === Packages ===
\usepackage[utf8]{inputenc}
\usepackage[T1]{fontenc}
\usepackage{amsmath,amssymb,amsthm}
\usepackage{mathtools}
\usepackage{algorithm}
\usepackage{algorithmic}
\usepackage{booktabs}
\usepackage{graphicx}
\usepackage{hyperref}
\usepackage{cleveref}
\usepackage{natbib}
\usepackage[margin=1in]{geometry}
\usepackage{xcolor}
% \usepackage{microtype} % Disabled: requires scalable fonts

% === Custom commands ===
\newcommand{\EU}{\mathrm{EU}}
\newcommand{\VOI}{\mathrm{VOI}}
\newcommand{\Beta}{\mathrm{Beta}}
\newcommand{\Prob}{\mathbb{P}}
\newcommand{\E}{\mathbb{E}}
\newcommand{\R}{\mathbb{R}}
\newcommand{\argmax}{\operatornamewithlimits{argmax}}
\newcommand{\credence}{\textsc{Credence}}

\newtheorem{definition}{Definition}
\newtheorem{proposition}{Proposition}
\newtheorem{axiom}{Axiom}

\title{Credence: A Bayesian Decision-Theoretic Framework \\for LLM Agent Tool Selection}

\author{
  Guy Freeman\thanks{Correspondence: \texttt{research@gfrm.in}. Blog: \url{https://gfrm.in}. Code: \url{https://github.com/gfrmin/credence}.} \\
  Independent Researcher, Tel Aviv
}

\date{\today}

\begin{document}
\maketitle

% ============================================================
% ABSTRACT
% ============================================================
\begin{abstract}
Current LLM agent frameworks select tools through prompting heuristics without any formal mechanism for evaluating whether a tool query is worth its cost. We introduce \credence{}, a Bayesian decision-theoretic framework that derives tool selection from first principles: agents maintain conjugate beliefs about tool reliability, update them from observed outcomes, and select queries by maximising expected utility via value-of-information calculations. The framework is grounded in five axioms---Bayesian conditioning, expected utility integration, EU maximisation, a Solomonoff complexity prior, and a cooperative inverse reinforcement learning (CIRL) alignment commitment---from which all agent behaviour is derived rather than engineered.

On a controlled benchmark, \credence{} scores $+163.7$ at zero API cost and $0.016$ seconds per question, dominating all non-LLM baselines by at least $3\times$. A frontier LLM (Claude Haiku~4.5) scores $+445.5$ with $97.5\%$ accuracy but costs \$3.24 per 1000 questions and runs $167\times$ slower; a local LLM (Llama~3.1 8B) scores $+79.3$---half the Bayesian agent---demonstrating that LLM-based tool selection without world knowledge underperforms principled statistical methods. Among non-LLM agents, we identify an \emph{accuracy paradox}: querying all tools achieves the highest accuracy ($64.4\%$) but the worst score ($-67.0$), because tool costs destroy the margin. Ablation confirms that reliability learning is the dominant component ($-159.9$ without it). These results establish \credence{} as the optimal zero-cost tool selection strategy and demonstrate that principled Bayesian methods offer a practical alternative when API cost or latency constraints preclude frontier model deployment.
\end{abstract}

% ============================================================
% 1. INTRODUCTION
% ============================================================
\section{Introduction}
\label{sec:intro}

The rapid deployment of large language model (LLM) agents---systems that use LLMs to select and invoke external tools---has created a curious gap in AI practice. While the perception capabilities of these systems have improved dramatically through scaling \citep{sutton2019bitter}, the \emph{decision-making} layer remains largely unprincipled. Frameworks such as LangChain \citep{chase2022langchain} and LlamaIndex implement tool selection through prompting: the LLM is instructed to reason about which tools to use, with no formal mechanism for quantifying uncertainty, evaluating marginal information value, or trading off query costs against expected gains.

This paper argues that the tool selection problem is fundamentally a \emph{decision-theoretic} problem, not a perception problem, and that treating it as such yields both theoretical clarity and practical improvements. Decision theory does not compete with deep learning for the same job---it answers a different question entirely: given uncertainty about tool reliability and finite query budgets, \emph{what should the agent do?}

We introduce \credence{}, a framework built on five axioms from which all agent behaviour is derived. The key insight is that tool queries are information-gathering actions whose value can be computed exactly: the \emph{value of information} (VOI) of querying a tool is the expected utility of observing the tool's response and then acting optimally, minus the expected utility of acting on current beliefs alone. A query is worth making if and only if its VOI exceeds its cost. This is not a heuristic---it is a theorem of Bayesian decision theory \citep{raiffa1961applied, degroot1984optimal}.

\paragraph{Contributions.} (1) We present a principled framework for LLM tool selection grounded in Bayesian decision theory and CIRL, requiring no LLM calls for the decision-making layer itself. (2) We identify the \emph{accuracy paradox}: among agents that depend on tool selection, higher accuracy can coexist with lower net value, because accuracy-maximising agents over-query. (3) We demonstrate empirically that a local LLM agent (Llama~3.1 8B) underperforms the Bayesian agent despite using native tool-calling, and that a frontier LLM (Claude Haiku~4.5) outperforms it only at non-trivial API cost (\$3.24 per 1000 questions) and $167\times$ higher latency. (4) Ablation confirms that reliability learning is the single most valuable component of the framework.

\paragraph{Connection to prior work.} The present framework builds on two lines of research. First, the theory of \emph{embodied agent design} due to \citet{ay2015geometric}, which formalises the separation between an agent's controller (brain) and its sensorimotor apparatus (body and environment), and argues that control architectures should be \emph{sufficient} for expressing desired behaviours while remaining \emph{concise}---the principle of cheap design. In \credence{}, the ``brain'' is the Bayesian decision layer; the ``body'' is the LLM and external tools. The LLM is treated as a \emph{prosthetic}---a noisy sensor whose reliability is learned, not a decision-maker.

Second, the author's prior work on \emph{dynamic staged trees} \citep{freeman2011dynamic} and \emph{chain event graph} (CEG) model selection \citep{freeman2011bayesian}, which demonstrated that Bayesian structure learning with conjugate updates can handle non-stationary discrete processes through multi-process switching. The Beta-Bernoulli reliability model in \credence{} is a specialisation of the Dirichlet-categorical conjugate framework used for CEG stage learning.


% ============================================================
% 2. BACKGROUND AND RELATED WORK
% ============================================================
\section{Related Work}
\label{sec:related}

\subsection{Principled Frameworks for LLM Agents}

The application of Bayesian decision theory to agent design has a long history \citep{russell1991principles, boutilier1999decision, kaelbling1998planning}, yet modern LLM agent frameworks have largely ignored these foundations. ReAct \citep{yao2023react}, Reflexion \citep{shinn2023reflexion}, and LATS \citep{zhou2024lats} delegate both perception and decision-making to the LLM (the latter via MCTS-style planning); Toolformer \citep{schick2023toolformer} fine-tunes models to decide when to call APIs. These approaches conflate two fundamentally different capabilities: extracting information from tool outputs (perception) and choosing which tools to invoke (decision-making). \citet{wang2023llmagent} survey over 100 such systems; none maintains explicit beliefs about tool reliability or computes value of information.

Three recent works apply principled methods to LLM agents. \textbf{RAFA} \citep{liu2024rafa} is the most formally rigorous: it casts agent reasoning as a Bayesian adaptive MDP and proves a $\sqrt{T}$ regret bound. However, RAFA addresses planning in unknown environments, not per-tool reliability or VOI for individual tool calls. \textbf{DeLLMa} \citep{liu2025dellma} applies expected utility maximisation to LLM decisions via multi-step inference---the closest in spirit to \credence{}'s EU framework---but lacks belief updating, VOI computation, and tool reliability tracking. \textbf{MACLA} \citep{forouzandeh2026macla} maintains Beta posteriors over procedure success rates, and its ablation confirms Bayesian selection as the single most impactful component. But MACLA selects multi-step procedures, not individual tools, and does not compute VOI.

To our knowledge, no published work combines per-tool Bayesian belief tracking, value-of-information computation, and expected utility maximisation for LLM tool selection.

\subsection{Cost-Aware Agent Approaches}

Cost-aware approaches to LLM tool use span five tiers of increasing formality: (1)~\emph{prompt-based}---expose budget or cost information in the prompt; (2)~\emph{RL reward shaping}---penalise tool calls in the reward signal; (3)~\emph{planning and optimisation}---pre-plan tool sequences under budget constraints, as in BTP \citep{zheng2024btp}; (4)~\emph{sequential decision-making}---online planning with cost anticipation, most rigorously formalised by INTENT \citep{liu2026intent}; and (5)~\emph{decision-theoretic}---computing VOI before each tool call, the approach taken by \credence{}.

\citet{kapoor2025agents} demonstrate that agent evaluation without cost accounting is misleading---accuracy alone fails to capture real-world value. \citet{wang2024tokeneconomies} show that budget-aware reasoning strategies outperform unconstrained ones, and \citet{erol2025costofpass} propose an economic framework for model evaluation that accounts for inference cost. \citet{wang2025epistemic} argue that agents should invoke tools ``only when epistemically necessary,'' defining a knowledge boundary concept but providing no computational mechanism. A Bayesian VOI framework provides exactly the mechanism they call for.

Our benchmark results (\S\ref{sec:results}) illustrate the cost-performance frontier: a local LLM with cost-awareness prompts and tool-calling still scores half the Bayesian agent, whilst a frontier LLM outperforms it only at non-trivial API cost. Natural language cannot express the quantitative tradeoff between information gain and query cost that VOI captures analytically. \citet{desabbata2024metareasoning} reach a complementary conclusion from the metareasoning perspective: LLMs benefit from explicit computational cost-benefit analysis rather than implicit prompt-level reasoning.

\subsection{Theoretical Foundations}

\paragraph{Embodied agent design.} \citet{ay2015geometric} formalises the sensorimotor loop as a causal model with Markov kernels and shows that an agent's controller need not be a universal approximator---it suffices to match the behavioural repertoire afforded by its embodiment. This \emph{cheap design} principle \citep{pfeifer2006body, montufar2015cheap} motivates \credence{}'s architecture: the Bayesian decision layer is far simpler than the LLM it coordinates, yet sufficient for optimal tool selection given the embodiment constraints (available tools and their reliability profiles).

\paragraph{Graphical models.} Chain event graphs \citep{smith2008chain, freeman2011bayesian, thwaites2009keod} generalise Bayesian networks to handle context-specific independence with conjugate model selection via product Dirichlet priors. Dynamic staged trees \citep{freeman2011dynamic} extend this to non-stationary time series through multi-process switching models. The per-tool, per-category Beta-Bernoulli reliability model in \credence{} is a specialisation of the product Dirichlet conjugate framework developed for chain event graph model selection \citep{freeman2011bayesian}, where \citet[Theorem~18]{freeman2011bayesian} showed that path and floret independence assumptions characterise the Dirichlet prior on staged tree parameters. Each tool-category pair in \credence{} corresponds to a floret in the CEG sense, with the correct/incorrect outcome defining a two-edge floret whose natural conjugate prior is Beta.

\paragraph{Value of information.} VOI originates with \citet{howard1966information} and is developed for bounded agents by \citet{russell1991principles}. Bayesian experimental design \citep{chaloner1995bayesian} formalises optimal information gathering under cost constraints. \citet{dong2026voi} recently extend VOI to human-agent communication, but apply it to dialogue acts rather than tool selection. \credence{} appears to be the first system to compute VOI for individual tool calls in an LLM agent.

\paragraph{Value alignment.} Cooperative inverse reinforcement learning \citep{hadfield2016cooperative} treats human-robot interaction as a cooperative game in which the robot infers human preferences from observed behaviour. \credence{}'s treatment of the LLM as a noisy sensor---whose outputs are observations with uncertainty, not commands---embodies a similar principle: the decision-theoretic layer maintains its own beliefs and defers to the human's utility function rather than to the LLM's token probabilities.


% ============================================================
% 3. THE CREDENCE FRAMEWORK
% ============================================================
\section{The \credence{} Framework}
\label{sec:framework}

\subsection{Axioms}
\label{sec:axioms}

\credence{} is built on five axioms. These are non-negotiable foundations from which all behaviour is derived.

\begin{axiom}[Conditioning]
Beliefs are updated by Bayes' rule. This is the unique update rule for which updating and computing the predictive commute, and the unique rule that cannot be Dutch-booked across time (diachronic coherence).
\end{axiom}

\begin{axiom}[Integration]
The expected value of a quantity under a belief is computed by integrating over the belief distribution.
\end{axiom}

\begin{axiom}[Optimisation]
Actions are selected to maximise expected utility under the current belief. Savage's representation theorem \citep{savage1954foundations} establishes that any agent whose preferences satisfy basic rationality axioms acts as if maximising expected utility. The complete class theorem \citep{wald1950statistical, berger1985statistical} further establishes that Bayesian decision rules are the only admissible ones: any non-Bayesian rule is dominated by some Bayesian rule.
\end{axiom}

\begin{axiom}[Complexity prior]
The prior over hypotheses is weighted by description length: $2^{-|h|}$ for hypothesis $h$. This is Solomonoff's universal prior \citep{solomonoff1964formal}---the unique prior that dominates all computable priors up to a constant.
\end{axiom}

\begin{axiom}[Alignment]
The agent's utility equals the user's utility under preference uncertainty: $U(a, s, b) = \E_{\theta \sim b}[u_\theta(s, a)]$, where $\theta$ parameterises the user's preferences and $b$ is the agent's belief over $\theta$. This is the cooperative inverse reinforcement learning (CIRL) framework of \citet{hadfield2016cooperative}: the agent and user are in a cooperative game where the agent's reward is the user's reward, but the agent does not know the user's reward function and must infer it from interaction.
\end{axiom}

\subsection{Architecture: Brain, Body, Environment}
\label{sec:architecture}

Following \citet{ay2015geometric}, we decompose the agent into three components:

\begin{itemize}
\item \textbf{Brain} (controller $\mathcal{C}$): The Bayesian decision layer. Maintains beliefs about tool reliability, computes expected utilities, selects actions. Implemented in ${\sim}200$ lines of Julia. Contains no LLM.
\item \textbf{Body} (actuators and sensors): The available tools, including the LLM itself. These are \emph{noisy channels} in the information-theoretic sense---Markov kernels $\beta: \mathcal{W} \to \Delta_\mathcal{S}$ from world states to sensor states, following the notation of \citet{ay2015geometric}. The LLM is a prosthetic: a powerful but unreliable sensor whose reliability profile is learned, not assumed.
\item \textbf{Environment} (world $\mathcal{W}$): The questions, their categories, and the ground-truth answers. The agent does not observe the environment directly---only through the noisy responses of its tools.
\end{itemize}

This separation instantiates the principle of cheap design: the brain is deliberately simple because the body (LLM + tools) provides the perceptual complexity. The brain's job is not to understand language or reason about content---it is to decide \emph{which tool to query and when to stop}.

\subsection{Belief Model}
\label{sec:beliefs}

For each tool $t \in \{1, \ldots, T\}$ and question category $c \in \{1, \ldots, C\}$, the agent maintains a Beta-distributed belief over the tool's reliability:
\begin{equation}
  p_{t,c} \sim \Beta(\alpha_{t,c}, \beta_{t,c}).
  \label{eq:beta}
\end{equation}

The conjugate Beta-Bernoulli model means that after observing a correct response from tool $t$ on a category-$c$ question, the update is:
\begin{equation}
  (\alpha_{t,c}, \beta_{t,c}) \mapsto (\alpha_{t,c} + 1, \beta_{t,c}),
\end{equation}
and after an incorrect response:
\begin{equation}
  (\alpha_{t,c}, \beta_{t,c}) \mapsto (\alpha_{t,c}, \beta_{t,c} + 1).
\end{equation}

\paragraph{Initialisation.} Priors are set to $\Beta(1, 1)$ (uniform) for each tool-category pair. This encodes maximum ignorance about reliability, placing the prior mean at $1/2$ with high variance. The choice of $1/N$ (where $N$ is the number of tools) for prior success probability, rather than $0.5$, was considered and rejected: the Beta(1,1) prior is the unique non-informative conjugate prior and makes no assumption about relative tool quality.

\paragraph{Non-stationarity.} The core framework handles distributional drift through standard Beta-Bernoulli updating: as a tool's reliability degrades, negative evidence accumulates and the posterior shifts downward naturally. No explicit change-detection mechanism is required---the posterior adapts continuously, with the rate of adaptation governed by the effective sample size (the sum $\alpha + \beta$).

\paragraph{Non-stationarity extensions.} The Beta-Bernoulli model adapts continuously to changing tool reliability through standard posterior updating: as negative evidence accumulates, the posterior shifts downward, governed by the effective sample size ($\alpha + \beta$). For abrupt regime changes, principled extensions include Bayesian online changepoint detection \citep{adams2007changepoint} and power steady models \citep{freeman2011dynamic}, which we defer to future work.

\subsection{Decision Rule: Value of Information}
\label{sec:voi}

The agent's decision at each step is: should I query another tool, or should I act (submit an answer or abstain) based on current beliefs?

Let $b$ denote the current belief state (the collection of all Beta parameters). Let $r_t$ denote the cost of querying tool $t$. Let $v_{\text{correct}} = 10$, $v_{\text{wrong}} = -5$, $v_{\text{abstain}} = 0$ be the payoffs.

\paragraph{Expected utility of acting now.}
Given current beliefs $b$, if the agent has gathered responses $\{(t_i, x_i)\}_{i=1}^k$ from $k$ tool queries, it computes a posterior probability $p_{\text{correct}}$ of the most-supported answer being correct (by aggregating the reliability-weighted evidence). The expected utility of submitting is:
\begin{equation}
  \EU_{\text{submit}}(b) = p_{\text{correct}} \cdot v_{\text{correct}} + (1 - p_{\text{correct}}) \cdot v_{\text{wrong}}.
\end{equation}
The expected utility of abstaining is $\EU_{\text{abstain}} = 0$. The agent submits if $\EU_{\text{submit}} > 0$, i.e., if $p_{\text{correct}} > 1/3$ (the break-even threshold given the payoff structure).

The best ``act now'' utility is:
\begin{equation}
  \EU^*(b) = \max(\EU_{\text{submit}}(b),\; 0).
\end{equation}

\paragraph{Value of information.}
The VOI of querying tool $t$ is:
\begin{equation}
  \VOI(t \mid b) = \E_{x \sim p_t}\!\bigl[\EU^*(b \mid x)\bigr] - \EU^*(b),
  \label{eq:voi}
\end{equation}
where the expectation is over the possible responses $x$ of tool $t$, weighted by the agent's current belief about tool $t$'s reliability. This is the expected improvement in decision quality from observing tool $t$'s response.

A query is made if and only if:
\begin{equation}
  \VOI(t \mid b) > r_t \quad \text{for some tool } t.
  \label{eq:query_criterion}
\end{equation}

When multiple tools satisfy this criterion, the agent selects the one with the highest net VOI: $\argmax_t [\VOI(t \mid b) - r_t]$.

\begin{proposition}
\label{prop:voi_decrease}
Under the Beta-Bernoulli model, $\VOI(t \mid b)$ is non-increasing in the number of prior observations. That is, the more the agent knows about a tool, the less informative an additional query becomes.
\end{proposition}

\begin{proof}
As $\alpha_{t,c} + \beta_{t,c} \to \infty$, the posterior variance $\mathrm{Var}(p_{t,c}) \to 0$ and the expected posterior after one additional observation converges to the current posterior. Hence $\E[\EU^*(b \mid x)] \to \EU^*(b)$ and $\VOI(t \mid b) \to 0$.
\end{proof}

This ensures that the agent's querying naturally terminates: after sufficiently many observations, no tool query is worth its cost.

\paragraph{Category-specific beliefs.}
The agent maintains separate Beta parameters for each tool-category pair. In the current benchmark, question categories are provided by the environment; the agent uses them to index into the appropriate reliability beliefs. This allows the agent to learn, for example, that the calculator is perfectly reliable for numerical questions but uninformative for others, without any text features or domain knowledge.


% ============================================================
% 4. EXPERIMENTAL SETUP
% ============================================================
\section{Experimental Setup}
\label{sec:setup}

\subsection{Simulated Environment}

We evaluate on a benchmark of 50 multiple-choice questions spanning five categories: factual, numerical, recent events, misconceptions, and reasoning. Tools are simulated with known reliability profiles (\cref{tab:tools}). While simulated tools limit ecological validity, they enable precise isolation of the decision-making layer's contribution and allow all agents to see identical tool responses for a given seed.

\begin{table}[t]
\centering
\caption{Tool reliability by question category. All tools respond to every question (no coverage mechanism). Reliability values represent $P(\text{correct} \mid \text{category})$; chance on 4-choice questions is $0.25$.}
\label{tab:tools}
\small
\begin{tabular}{@{}lcccccr@{}}
\toprule
& \multicolumn{5}{c}{$P(\text{correct})$} & Cost \\
\cmidrule(lr){2-6}
Tool & Fact. & Num. & Rec. & Misc. & Reas. & \\
\midrule
A (Web Search)      & 0.70 & 0.20 & 0.65 & 0.25 & 0.40 & 1 \\
B (Knowledge Base)  & 0.92 & 0.40 & 0.55 & 0.88 & 0.45 & 2 \\
C (Calculator)      & 0.25 & 1.00 & 0.25 & 0.25 & 0.25 & 1 \\
D (LLM Direct)      & 0.65 & 0.50 & 0.45 & 0.40 & 0.72 & 2 \\
\bottomrule
\end{tabular}
\end{table}

\paragraph{Scoring.} Correct answer: $+10$. Wrong answer: $-5$. Abstention: $0$. Tool costs as in \cref{tab:tools}. The asymmetric payoff creates a break-even submission threshold of $p_{\text{correct}} > 1/3$.

\subsection{Agents Compared}

\paragraph{\credence{} (Bayesian).} The full framework of \cref{sec:framework}. Beta(1,1) priors per tool-category pair, VOI-based tool selection. No LLM calls. Implemented in ${\sim}200$ lines of Julia.

\paragraph{LLM agents.} Two LLM-based agents, both using native tool-calling (structured API, not text parsing). Both receive the same system prompt describing the four tools, scoring rules, the option to abstain, and a rolling history of the last 10 outcomes.
\begin{itemize}
\item \textbf{Claude Haiku 4.5} (claude-haiku-4-5-20251001): Anthropic's frontier model. API cost: \$1/M input, \$5/M output tokens. Represents the upper bound of LLM-based tool selection.
\item \textbf{Llama 3.1 8B}: local model via Ollama, temperature 0. Zero API cost. Represents the zero-cost LLM baseline.
\end{itemize}
The comparison is between decision-making paradigms---Bayesian EU maximisation vs.\ LLM-prompted reasoning---not software frameworks.

\paragraph{Single-Best-Tool.} Always queries Tool~A (web search, cost~1) and submits its answer. Tests whether \credence{}'s advantage is merely tool routing.

\paragraph{All-Tools.} Queries all four tools, majority-votes. Demonstrates the cost of indiscriminate querying.

\paragraph{Random.} Queries a random tool, submits its response. Lower bound.

\subsection{Experimental Protocol}

All results are averaged over 20 random seeds, with 50 questions per seed. Tool responses are pre-generated per seed so that all agents see identical (question, tool) outputs, ensuring a fair comparison. LLM agents use native tool-calling via structured API. Tool reliabilities are fixed as in \cref{tab:tools} throughout.


% ============================================================
% 5. RESULTS
% ============================================================
\section{Results}
\label{sec:results}

\subsection{Stationary Environment}

\begin{table}[t]
\centering
\caption{Benchmark results (mean $\pm$ std over 20 seeds, 50 questions per seed). API Cost is the total across all 1000 questions.}
\label{tab:stationary}
\small
\begin{tabular}{@{}lrrrrrrr@{}}
\toprule
Agent & Score & Acc. & Abst.\% & Tools/Q & Cost/Q & Time/Q (s) & API Cost \\
\midrule
Claude Haiku 4.5 & $\mathbf{+445.5 \pm 13.9}$ & \textbf{0.975} & 0.000 & 0.59 & 0.71 & 2.67 & \$3.24 \\
\credence{} (Bayesian) & $\mathbf{+163.7 \pm 51.7}$ & 0.647 & 0.045 & 0.96 & 1.25 & 0.016 & \$0.00 \\
Llama 3.1 8B & $+79.3 \pm 36.7$ & 0.594 & 0.229 & 0.90 & 1.44 & 1.56 & \$0.00 \\
Random & $+55.4 \pm 48.7$ & 0.509 & 0.000 & 1.00 & 1.53 & ${\sim}0$ & \$0.00 \\
Single-Best-Tool & $+44.2 \pm 50.1$ & 0.459 & 0.000 & 1.00 & 1.00 & ${\sim}0$ & \$0.00 \\
All-Tools & $-67.0 \pm 65.7$ & 0.644 & 0.000 & 4.00 & 6.00 & ${\sim}0$ & \$0.00 \\
\bottomrule
\end{tabular}
\end{table}


\Cref{tab:stationary} reveals a cost-performance frontier. Claude Haiku~4.5 dominates on raw score ($+445.5$), answering $97.5\%$ of questions correctly whilst using fewer tools per question ($0.59$) than any other agent---it answers over half of all questions from world knowledge alone, without querying any tool. However, this performance costs \$3.24 per 1000 questions and $2.67$ seconds per question. \credence{} scores $+163.7$ at zero API cost and $0.016$ seconds per question---$37\%$ of the frontier LLM's score at $0\%$ of its cost and $0.6\%$ of its latency.

Haiku~4.5's high score is partly attributable to answering $54\%$ of questions from world knowledge without tool use---a capability the Bayesian agent lacks by design, since it has no internal knowledge and \emph{must} query tools. The comparison is most informative among agents that depend on tool selection for their performance.

\paragraph{The accuracy paradox.}
Among non-LLM agents, higher accuracy coexists with lower net value. All-Tools achieves the highest non-LLM accuracy ($64.4\%$) but the worst score ($-67.0$) because querying all four tools costs $6.0$ per question, destroying the margin. \credence{} achieves comparable accuracy ($64.7\%$) whilst spending only $1.25$ per question on tools. The net score is:
\begin{equation}
  \text{Score} = \sum_{q=1}^Q \bigl[v_q \cdot \mathbf{1}[\text{correct}_q] + v_{\text{wrong}} \cdot \mathbf{1}[\text{wrong}_q]\bigr] - \sum_{q=1}^Q \text{cost}_q,
\end{equation}
where $v_q = v_{\text{correct}}$ if the agent submits correctly, $v_{\text{wrong}}$ if incorrectly, and $0$ if abstaining. Maximising the first sum (accuracy) is not equivalent to maximising the full expression when costs are non-zero.

\paragraph{Bayesian vs.\ local LLM.}
Llama~3.1 8B scores $+79.3$---less than half the Bayesian agent's score---despite using native tool-calling. It abstains too aggressively ($22.9\%$ vs.\ $4.5\%$) and achieves lower accuracy on submitted answers ($59.4\%$ vs.\ $64.7\%$). A local 8B model lacks the world knowledge to compete with either the Bayesian agent's principled tool selection or a frontier model's parametric knowledge.

\paragraph{Bayesian vs.\ baselines.}
Among zero-cost agents, \credence{} dominates: it scores $3.0\times$ higher than Random ($+55.4$) and $3.7\times$ higher than Single-Best-Tool ($+44.2$). The gap reflects the combined value of VOI-based tool selection, reliability learning, and selective abstention.

\subsection{Ablation Study}

\begin{table}[t]
\centering
\caption{Ablation study: effect of removing individual components (20 seeds, 50 questions per seed).}
\label{tab:ablation}
\begin{tabular}{@{}lrrrl@{}}
\toprule
Variant & Score & $\Delta$ & Tools/Q & Effect \\
\midrule
Full \credence{} & $+163.7 \pm 51.7$ & --- & 0.96 \\
Greedy best & $+189.4 \pm 38.9$ & $+25.7$ & 1.00 & Picks most reliable tool; no VOI needed \\
No abstention & $+156.9 \pm 60.4$ & $-6.8$ & 0.96 & Forced submissions on low confidence \\
No VOI (query all) & $+38.2 \pm 56.8$ & $-125.5$ & 4.00 & Highest accuracy (79.1\%) but cost destroys margin \\
No learning & $+3.8 \pm 44.7$ & $-159.9$ & 1.00 & Barely above zero; cannot identify good tools \\
\bottomrule
\end{tabular}
\end{table}

Reliability learning is the dominant component: removing it ($-159.9$) reduces the agent to near-zero performance, because without learning the agent cannot identify which tools are reliable for which question categories. Removing VOI ($-125.5$) forces the agent to query all four tools on every question; it achieves the highest accuracy of any non-LLM agent ($79.1\%$) but the cost of $6.0$ per question destroys the margin---the same mechanism that makes All-Tools the worst-scoring agent in \cref{tab:stationary}. Abstention contributes modestly ($-6.8$).

The greedy variant---which queries the single tool with highest expected reliability and submits its answer directly---slightly outperforms the full framework ($+189.4$ vs.\ $+163.7$). This occurs because the greedy agent commits to the most reliable tool per category (often the more expensive Knowledge Base at cost~$2$), whilst the full agent's VOI calculation sometimes favours cheaper but less reliable tools. The greedy agent also never abstains, and in this payoff structure the forced submissions are net positive. However, the greedy agent's success depends entirely on learning: without it, greedy reduces to picking an arbitrary tool. The key insight is that \emph{learning which tools are reliable is the primary contribution}; VOI provides the mechanism to prevent wasteful over-querying when tool costs are significant relative to information gain.


% ============================================================
% 6. DISCUSSION
% ============================================================
\section{Discussion}
\label{sec:discussion}

\paragraph{Why accuracy is the wrong metric.}
The accuracy paradox is not merely an artefact of our scoring system. Any deployment context with non-zero query costs---API fees, latency, rate limits---creates a divergence between accuracy and net value. Among non-LLM agents, All-Tools achieves the highest accuracy ($64.4\%$) but the worst score ($-67.0$). Even for the frontier LLM, the relevant metric is score-per-dollar: Haiku~4.5 achieves $2.7\times$ the Bayesian agent's score but at infinite cost relative to the Bayesian agent's zero. The correct metric is expected utility, which accounts for both the value of correct answers and the cost of obtaining them.

\paragraph{Cheap design for LLM agents.}
\citet{ay2015geometric} argues that an embodied agent's controller need not be a universal approximator---it need only be sufficient for expressing desired behaviours given the agent's embodiment constraints. The \credence{} decision layer demonstrates this principle: a Julia inference engine maintaining Beta distributions and computing VOI outperforms a local LLM-based decision-maker that is orders of magnitude more complex, and achieves $37\%$ of a frontier model's score at zero cost. The LLM is powerful but expensive; using it for decisions that admit analytical solutions---resource allocation under uncertainty---is the opposite of cheap design.

The analogy to Ay's exponential family policy models (\citealt{ay2015geometric}, Proposition~1) is suggestive: just as Ay shows that policies parameterised by a small number of ``feature vectors'' derived from the embodiment constraints can be embodied universal approximators, the \credence{} decision layer achieves strong performance using only the sufficient statistics of the Beta posteriors. The ``feature vectors'' are the tool-category reliability estimates; the ``embodiment constraints'' are the tools' true (unknown) reliability profiles.

The computational asymmetry is stark: the \credence{} decision layer processes each question in $0.016$ seconds (closed-form Beta posterior updates and VOI computation), whilst the frontier LLM requires $2.67$ seconds---a $167\times$ difference driven entirely by LLM inference overhead for tool selection decisions that admit analytical solutions. \Cref{tab:stationary} reports wall-clock time per question for each agent.

\paragraph{Connection to dynamic staged trees.}
The per-tool, per-category Beta parameters in \credence{} are structurally analogous to the stage parameters in a chain event graph \citep{freeman2011bayesian}: each tool-category pair is a ``situation'' in the event tree sense, with its own Dirichlet-distributed transition probabilities. The agent's task of learning which tools are reliable for which categories mirrors the CEG model selection problem of learning which situations share the same stage (i.e., the same transition distribution).

\paragraph{Limitations.}
The current evaluation uses simulated tools with known ground truth, which limits ecological validity. The question bank is small (50 questions) and question categories are provided to the agent rather than inferred. We include a frontier LLM (Claude Haiku~4.5), which comprehensively outperforms \credence{} on raw score ($+445.5$ vs.\ $+163.7$), demonstrating that world knowledge accessible through a frontier model's parameters can compensate for the lack of principled tool selection. However, this comes at an API cost (\$3.24 per 1000 questions) and latency ($2.67$s per question) that may be prohibitive at scale. The current framework assumes tools return independent responses---correlated errors between tools are not modelled. Distributional drift experiments are deferred to future work.


% ============================================================
% 7. CONCLUSION AND FUTURE WORK
% ============================================================
\section{Conclusion and Future Work}
\label{sec:conclusion}

We have presented \credence{}, a Bayesian decision-theoretic framework for LLM tool selection that derives agent behaviour from five axioms rather than engineering it through prompting. On a controlled benchmark, \credence{} achieves the highest score among zero-cost agents ($+163.7$), dominating all non-LLM baselines by at least $3\times$. A frontier LLM (Claude Haiku~4.5) outperforms it on raw score ($+445.5$) but at non-trivial API cost and $167\times$ higher latency. These results establish a cost-performance frontier: principled Bayesian methods are optimal when cost or latency constraints preclude frontier model deployment, and remain relevant as a decision-making layer even when frontier models are available.

Several extensions are immediate. First, evaluation with real (non-simulated) tools to establish ecological validity. Second, distributional drift experiments to evaluate the framework's robustness to non-stationary tool reliability. Third, the full \credence{} architecture extends beyond tool selection to a general-purpose agent framework with programs represented as DSL expressions (closed-loop policies in the sense of \citealt{sutton1999between}), meta-actions that modify the hypothesis space, and grammar evolution for discovering abstractions---a ``strange loop'' where the agent improves its own capacity for improvement. This broader architecture, grounded in the same five axioms, will be presented in a companion paper.

Third, the CIRL alignment axiom (\cref{sec:axioms}) opens a natural connection to the preference learning literature: the agent's beliefs about user preferences can be updated from implicit feedback (clicks, corrections, task completions) using the same conjugate machinery that tracks tool reliability. This positions \credence{} not merely as an agent framework but as an alignment-by-construction approach, where safety properties follow from the axioms rather than being imposed as constraints.

\paragraph{Reproducibility.} Code and data are available at \url{https://github.com/gfrmin/credence}. All experiments use fixed random seeds.


% ============================================================
% BIBLIOGRAPHY
% ============================================================
\bibliographystyle{plainnat}

\begin{thebibliography}{99}

\bibitem[Adams and MacKay(2007)]{adams2007changepoint}
R.~P. Adams and D.~J.~C. MacKay.
\newblock Bayesian online changepoint detection.
\newblock arXiv preprint arXiv:0710.3742, 2007.

\bibitem[Ay(2015)]{ay2015geometric}
N.~Ay.
\newblock Geometric design principles for brains of embodied agents.
\newblock \emph{K{\"u}nstliche Intelligenz}, 29(4):389--399, 2015.

\bibitem[Berger(1985)]{berger1985statistical}
J.~O. Berger.
\newblock \emph{Statistical Decision Theory and Bayesian Analysis}.
\newblock Springer, 2nd edition, 1985.

\bibitem[Boutilier et~al.(1999)]{boutilier1999decision}
C.~Boutilier, T.~Dean, and S.~Hanks.
\newblock Decision-theoretic planning: Structural assumptions and computational leverage.
\newblock \emph{J. Artificial Intelligence Research}, 11:1--94, 1999.

\bibitem[Chaloner and Verdinelli(1995)]{chaloner1995bayesian}
K.~Chaloner and I.~Verdinelli.
\newblock Bayesian experimental design: A review.
\newblock \emph{Statistical Science}, 10(3):273--304, 1995.

\bibitem[Chase(2022)]{chase2022langchain}
H.~Chase.
\newblock {LangChain}.
\newblock \url{https://github.com/langchain-ai/langchain}, 2022.

\bibitem[Chen et~al.(2023)]{chen2023chatgpt}
L.~Chen, M.~Zaharia, and J.~Zou.
\newblock How is {ChatGPT}'s behavior changing over time?
\newblock arXiv preprint arXiv:2307.09009, 2023.

\bibitem[De~Sabbata et~al.(2024)]{desabbata2024metareasoning}
C.~N. De~Sabbata, T.~R. Sumers, B.~AlKhamissi, A.~Bosselut, and T.~L. Griffiths.
\newblock Rational metareasoning for large language models.
\newblock In \emph{NeurIPS 2024 Workshop}, 2024.

\bibitem[DeGroot(1984)]{degroot1984optimal}
M.~H. DeGroot.
\newblock \emph{Optimal Statistical Decisions}.
\newblock John Wiley \& Sons, 1984.

\bibitem[Dong et~al.(2026)]{dong2026voi}
Y.~R. Dong, T.~Hu, Z.~Hui, C.~Zhang, I.~Vuli\'{c}, A.~Bobu, and N.~Collier.
\newblock Value of information: A framework for human-agent communication.
\newblock arXiv preprint arXiv:2601.06407, 2026.

\bibitem[Erol et~al.(2025)]{erol2025costofpass}
M.~H. Erol, B.~El, M.~Suzgun, M.~Yuksekgonul, and J.~Zou.
\newblock Cost-of-pass: An economic framework for evaluating language models.
\newblock arXiv preprint arXiv:2504.13359, 2025.

\bibitem[Forouzandeh et~al.(2026)]{forouzandeh2026macla}
S.~Forouzandeh, W.~Peng, P.~Moradi, X.~Yu, and M.~Jalili.
\newblock Learning hierarchical procedural memory for {LLM} agents through {B}ayesian selection and contrastive refinement.
\newblock In \emph{International Conference on Autonomous Agents and Multi-Agent Systems (AAMAS)}, 2026.

\bibitem[Freeman and Smith(2011a)]{freeman2011bayesian}
G.~Freeman and J.~Q. Smith.
\newblock Bayesian {MAP} model selection of chain event graphs.
\newblock \emph{J. Multivariate Analysis}, 102(7):1152--1165, 2011.

\bibitem[Freeman and Smith(2011b)]{freeman2011dynamic}
G.~Freeman and J.~Q. Smith.
\newblock Dynamic staged trees for discrete multivariate time series: Forecasting, model selection and causal analysis.
\newblock \emph{Bayesian Analysis}, 6(2):279--305, 2011.

\bibitem[Hadfield-Menell et~al.(2016)]{hadfield2016cooperative}
D.~Hadfield-Menell, S.~J. Russell, P.~Abbeel, and A.~Dragan.
\newblock Cooperative inverse reinforcement learning.
\newblock In \emph{Advances in Neural Information Processing Systems}, 2016.

\bibitem[Howard(1966)]{howard1966information}
R.~A. Howard.
\newblock Information value theory.
\newblock \emph{IEEE Trans.\ Systems Science and Cybernetics}, 2(1):22--26, 1966.

\bibitem[Kaelbling et~al.(1998)]{kaelbling1998planning}
L.~P. Kaelbling, M.~L. Littman, and A.~R. Cassandra.
\newblock Planning and acting in partially observable stochastic domains.
\newblock \emph{Artificial Intelligence}, 101(1-2):99--134, 1998.

\bibitem[Kapoor et~al.(2025)]{kapoor2025agents}
S.~Kapoor, B.~Stroebl, Z.~S. Siegel, N.~Nadgir, and A.~Narayanan.
\newblock {AI} agents that matter.
\newblock \emph{Transactions on Machine Learning Research}, 2025.

\bibitem[Liu et~al.(2024)]{liu2024rafa}
Z.~Liu, H.~Hu, S.~Zhang, H.~Guo, S.~Ke, B.~Liu, and Z.~Wang.
\newblock Reason for future, act for now: A principled framework for autonomous {LLM} agents with provable sample efficiency.
\newblock In \emph{International Conference on Machine Learning (ICML)}, 2024.

\bibitem[Liu et~al.(2025)]{liu2025dellma}
O.~Liu, D.~Fu, D.~Yogatama, and W.~Neiswanger.
\newblock {DeLLMa}: Decision making under uncertainty with large language models.
\newblock In \emph{International Conference on Learning Representations (ICLR)}, 2025.

\bibitem[Liu et~al.(2026)]{liu2026intent}
H.~Liu, C.~Tian, N.~An, Z.~Wang, P.~Lu, C.~Yu, and Q.~Qi.
\newblock Budget-constrained agentic large language models: Intention-based planning for costly tool use.
\newblock arXiv preprint arXiv:2602.11541, 2026.

\bibitem[Mont\'{u}far et~al.(2015)]{montufar2015cheap}
G.~Mont\'{u}far, K.~Ghazi-Zahedi, and N.~Ay.
\newblock A theory of cheap control in embodied systems.
\newblock \emph{PLOS Computational Biology}, 11(9):e1004427, 2015.

\bibitem[Pfeifer and Bongard(2006)]{pfeifer2006body}
R.~Pfeifer and J.~C. Bongard.
\newblock \emph{How the Body Shapes the Way We Think}.
\newblock MIT Press, 2006.

\bibitem[Raiffa and Schlaifer(1961)]{raiffa1961applied}
H.~Raiffa and R.~Schlaifer.
\newblock \emph{Applied Statistical Decision Theory}.
\newblock Harvard University Press, 1961.

\bibitem[Russell and Wefald(1991)]{russell1991principles}
S.~Russell and E.~Wefald.
\newblock \emph{Do the Right Thing: Studies in Limited Rationality}.
\newblock MIT Press, 1991.

\bibitem[Savage(1954)]{savage1954foundations}
L.~J. Savage.
\newblock \emph{The Foundations of Statistics}.
\newblock John Wiley \& Sons, 1954.

\bibitem[Schick et~al.(2023)]{schick2023toolformer}
T.~Schick, J.~Dwivedi-Yu, R.~Dess\`{i}, R.~Raileanu, M.~Lomeli, L.~Zettlemoyer, N.~Cancedda, and T.~Scialom.
\newblock Toolformer: Language models can teach themselves to use tools.
\newblock In \emph{Advances in Neural Information Processing Systems}, 2023.

\bibitem[Shinn et~al.(2023)]{shinn2023reflexion}
N.~Shinn, F.~Cassano, E.~Berman, A.~Gopinath, K.~Narasimhan, and S.~Yao.
\newblock Reflexion: Language agents with verbal reinforcement learning.
\newblock In \emph{Advances in Neural Information Processing Systems}, 2023.

\bibitem[Smith and Anderson(2008)]{smith2008chain}
J.~Q. Smith and P.~E. Anderson.
\newblock Conditional independence and chain event graphs.
\newblock \emph{Artificial Intelligence}, 172(1):42--68, 2008.

\bibitem[Solomonoff(1964)]{solomonoff1964formal}
R.~J. Solomonoff.
\newblock A formal theory of inductive inference.
\newblock \emph{Information and Control}, 7(1):1--22, 1964.

\bibitem[Sutton(2019)]{sutton2019bitter}
R.~S. Sutton.
\newblock The bitter lesson.
\newblock \url{http://www.incompleteideas.net/IncIdeas/BitterLesson.html}, 2019.

\bibitem[Sutton et~al.(1999)]{sutton1999between}
R.~S. Sutton, D.~Precup, and S.~Singh.
\newblock Between {MDPs} and semi-{MDPs}: A framework for temporal abstraction in reinforcement learning.
\newblock \emph{Artificial Intelligence}, 112(1-2):181--211, 1999.

\bibitem[Thwaites et~al.(2009)]{thwaites2009keod}
P.~A. Thwaites, G.~Freeman, and J.~Q. Smith.
\newblock Chain event graph {MAP} model selection.
\newblock In \emph{Proceedings of KEOD 2009}, pages 392--395, 2009.

\bibitem[Wald(1950)]{wald1950statistical}
A.~Wald.
\newblock \emph{Statistical Decision Functions}.
\newblock John Wiley \& Sons, 1950.

\bibitem[Wang et~al.(2023)]{wang2023llmagent}
L.~Wang, C.~Ma, X.~Feng, Z.~Zhang, H.~Yang, J.~Zhang, Z.~Chen, J.~Tang, X.~Chen, Y.~Lin, W.~X. Zhao, Z.~Wei, and J.-R. Wen.
\newblock A survey on large language model based autonomous agents.
\newblock \emph{Frontiers of Computer Science}, 18(6):186345, 2024.

\bibitem[Wang et~al.(2024)]{wang2024tokeneconomies}
J.~Wang, S.~Jain, D.~Zhang, B.~Ray, V.~Kumar, and B.~Athiwaratkun.
\newblock Reasoning in token economies: Budget-aware evaluation of {LLM} reasoning strategies.
\newblock In \emph{Conference on Empirical Methods in Natural Language Processing (EMNLP)}, 2024.

\bibitem[Wang et~al.(2025)]{wang2025epistemic}
H.~Wang, C.~Qian, M.~Li, J.~Qiu, B.~Xue, M.~Wang, H.~Ji, A.~Storkey, and K.-F. Wong.
\newblock Position: Agent should invoke external tools only when epistemically necessary.
\newblock In \emph{International Conference on Machine Learning (ICML)}, 2025.

\bibitem[Yao et~al.(2023)]{yao2023react}
S.~Yao, J.~Zhao, D.~Yu, N.~Du, I.~Shafran, K.~Narasimhan, and Y.~Cao.
\newblock {ReAct}: Synergizing reasoning and acting in language models.
\newblock In \emph{International Conference on Learning Representations}, 2023.

\bibitem[Zheng et~al.(2024)]{zheng2024btp}
Y.~Zheng, P.~Li, M.~Yan, J.~Zhang, F.~Huang, and Y.~Liu.
\newblock Budget-constrained tool learning with planning.
\newblock In \emph{Findings of the Association for Computational Linguistics (ACL)}, 2024.

\bibitem[Zhou et~al.(2024)]{zhou2024lats}
A.~Zhou, K.~Yan, M.~Shlapentokh-Rothman, H.~Wang, and Y.-X. Wang.
\newblock Language agent tree search unifies reasoning acting and planning in language models.
\newblock In \emph{International Conference on Machine Learning (ICML)}, 2024.

\end{thebibliography}

\end{document}
